\runninghead{Scope and Qualifications}
\section{Appendix: Scope and Qualifications}

\subsection{Standing Notice on Rights and Study English}
The approved proclaimed liturgical texts of the United States --- \work{The Roman Missal}, Third Typical Edition (U.S.\ English, copyright ICEL, 2010/2011), the protected modern Latin texts of composed Missal orations, Prefaces, and solemn blessings, and the \work{Lectionary for Mass for Use in the Dioceses of the United States of America} (New American Bible and Revised Grail / Abbey Psalms) --- are protected works and are neither reproduced nor closely paraphrased. All English Scripture quotations follow the tracked public-domain Douay--Rheims (Challoner) translation strictly as \textbf{public-domain study English}, never as the approved proclaimed English of the U.S.\ Lectionary. All Latin Scripture texts follow the public-domain Clementine Vulgate; the Latin Communion Antiphon (\emph{Dominus dabit benignitatem, et terra nostra dabit fructum suum}) is additionally collated against the secondary 2002 digital reproduction of the third typical edition of the \work{Missale Romanum} (p.~83).

\subsection{Liturgical Compilation and Interpretive Synthesis}
The postconciliar First Sunday of Advent in Year~B combines Missal prayers and antiphons shared across Years~A, B, and~C with a Year~B course of readings arranged in the \work{Ordo lectionum Missae}. The patristic and medieval commentaries cited in this study expound individual biblical passages or the Advent season as their authors received them; bringing compatible patristic arguments together across all twelve parts of this Sunday's Mass is an editorial synthesis grounded in those checked loci, not a claim that any Church Father preached on this exact postconciliar combination as one pre-existing sermon.

\subsection{Vulgate and Hebrew Numbering}
Because the restricted U.S.\ Lectionary day page and 2010 ICEL Antiphonary were unavailable offline, the chronological table in the \hyperref[sec:date-location]{Scriptural Date and Location} appendix follows the Vulgate numbering of the inspected Latin liturgical controls: Isaiah~63:16b--17; 64:1, 3b--8 (corresponding to U.S.\ Lectionary Hebrew/NABRE numbering Isa.~63:16b--17, 19b; 64:2b--7) and Vulgate Psalms~24, 79, and~84 (corresponding to Hebrew/Lectionary Psalms~25, 80, and~85). Fractional liturgical excerpts are queried in the chronological corpus at whole-verse canonical granularity; no subverse chronology or appointed modern wording is inferred from those envelopes.

\subsection{Verification and Offline Boundaries}
All primary research was conducted offline against held texts and verified records without reopening restricted remote payloads. Accordingly, the following points remain explicit open verification boundaries rather than asserted facts:
\begin{enumerate}
\item \textbf{Printed 2008 \work{Missale Romanum} correspondence:} Printed altar-book pages of the 2008 \work{reimpressio emendata} for \work{Dominica I Adventus} and the Advent entries of \work{Notitiae}~44 (2008), nn.~503--504 (pp.~368--387) were not reopened offline; collation of the Communion Antiphon and Prayer over the Offerings rests on verified records against the secondary 2002 digital PDF, leaving printed 2002 and printed 2008 correspondence unverified.
\item \textbf{Entrance Antiphon cue and truncation:} The biblical reference Psalm~25(24):1--3 rests on a cataloged lead; exact Missal cue wording and editorial truncation were not directly collated from the 2002/2008 Missal or 2010 ICEL Antiphonary.
\item \textbf{Modern Missal orations, Advent Preface details, and optional Solemn Blessings:} Modern Latin and ICEL English bodies for \emph{Da, quaesumus, omnipotens Deus}, \emph{Suscipe, quaesumus, Domine}, and \emph{Prosint nobis}, the exact Latin subtitle, date-range rubric, and text of the First Sunday of Advent Preface (beyond the verified count of two Prefaces of Advent in \work{Ordo Missae, Praefationes}), and the specific Advent solemn blessing numbering, title, and wording were not inspected offline and are neither quoted nor reconstructed.
\item \textbf{Calendar day-entry, U.S.\ Common Psalms, Appendix Eucharistic Prayers, and unextracted GIRM articles:} The individual November~29 day-entry page of the restricted USCCB 2027 calendar PDF, the live USCCB daily readings web page, the exact copyrighted refrains and verse spans of the two U.S.\ Advent Common Psalms under Lectionary No.~174 (Ps.~25[24] and Ps.~85[84]), the seasonal compatibility of the Appendix Eucharistic Prayers (Reconciliation; Various Needs and Occasions), and unextracted GIRM articles (GIRM~54 on the Collect, GIRM~63's unsung-omission clause, GIRM~305 and~313 on Advent altar and musical moderation, and GIRM~346 on liturgical color) remain unverified offline.
\end{enumerate}

\subsection{Textual Quality of Historical and Patristic Copies}
Historical sacramentary prayers (\work{Sacramentarium Gelasianum Vetus} II.lxxxiv and \work{Sacramentarium Veronense} XXVIIII.372) are cited strictly as patristic-era Roman antecedents and comparative controls, never as substitutes for the modern Missal prayers. Three specific textual-copy qualifications apply to the commentaries cited in the exposition:
\begin{itemize}
\item \textbf{St.~Bede, \work{In Marci Evangelium Expositio} IV.13:} All Latin wording and doctrinal summaries follow the inspected Wikisource and Corpus Corporum digital transcription without claiming print-facsimile collation against Migne's \work{Patrologia Latina} Vol.~92.
\item \textbf{Uncollated OCR locating copies:} Citations from St.~Jerome (\work{Commentarii in Isaiam} XVII, PL~24 cols.~642--651), St.~Thomas Aquinas (\work{Expositio super Isaiam ad litteram} ch.~63--64, Parma Vol.~14, pp.~370--371), Cornelius a Lapide (\work{In 1~Cor.}~1:3--4, 6, 8--9, Vol.~7; whereas Lapide on 1~Cor.~1:5 and~1:7 rests on facsimile-checked passages), and \work{Sacramentarium Gelasianum Vetus} II.lxxxiv (Wilson 1894, p.~218, preserving one unreadable OCR token before \emph{voluntatem}) rest on uncollated optical character recognition copies. In PL~24, Vallarsi's editorial footnote suggesting Origen is kept strictly separate from Jerome's own text.
\item \textbf{Mediated \work{Catena Aurea} citations:} St.~Augustine (\work{Epistula}~199 §3), Theophylact of Ohrid, and Pseudo-Jerome on Mark~13:32--37 are cited explicitly as mediated through St.~Thomas Aquinas's \work{Catena Aurea in Marcum} (1842 Oxford English edition).
\end{itemize}

\subsection{Canonical Audit Records}
Full Missal-level collation records, edition locators, and ancient sacramentary audits are maintained exclusively in the canonical shared Missal owner for the First Sunday of Advent (\texttt{PC-S01}); cycle-specific Lectionary audits, occurrence controls, source dispositions, and detailed reception matrices live in the companion research and instance records accompanying this publication.
